\section{Materialize}

\paragraph{Input and output.}
Materialization consumes a generic one-file Axon AST and checkpoint metadata.
It emits a one-file Axon AST specialized enough to be inspected or used as a
checkpoint-specific artifact.

\paragraph{Contract.}
The stage may resolve materialization-specific configuration expressions, but
it must not become an optimizer.  It does not perform dead-code elimination,
general constant folding, alias expansion, inlining, type inference, or
backend-specific rewriting.

\paragraph{Rationale.}
Generic model files are the source of truth for model families.  Materialized
files are useful for inspection and targeted debugging, but materialization
must not silently repair or reinterpret the program in ways that diverge from
the generic source.

\paragraph{Formal view.}
For generic AST \(A_f\) and checkpoint metadata \(C\), materialization is a
source-to-source transformation
\[
  \mathsf{materialize}(A_f,C) = A_f'
\]
that may replace materialization-approved configuration expressions with
checkpoint-specific values.  It must preserve the program's definition graph
and must not depend on backend selection.

\paragraph{Example.}
A generic constant may be written as
\begin{verbatim}
MODEL_DIM = Config.int@@hidden_size default=768
\end{verbatim}
For a checkpoint whose configuration contains \verb|hidden_size = 1024|,
materialization may emit
\begin{verbatim}
MODEL_DIM = 1024
\end{verbatim}
It should not also inline all uses of \verb|MODEL_DIM|, delete unused branches,
or specialize unrelated functions.
